% Unidad: métrica inducida por una norma (castellano).

\begin{frame}
\didactatitle{La métrica inducida}

Toda norma induce una métrica.

\begin{proposition}[Métrica inducida]
Si $(E, \|\cdot\|)$ es un espacio normado, entonces
\[
  d(x,y) = \|x - y\|
\]
define una métrica en $E$.
\end{proposition}

\onlynotes{%
\begin{proof}
Las tres propiedades de métrica se siguen directamente de los axiomas de norma.
La positividad y la anulación vienen de la positividad de la norma. La simetría
sale de la homogeneidad con $\lambda = -1$:
\[
  d(y,x) = \|y-x\| = \|-(x-y)\| = |-1|\,\|x-y\| = d(x,y).
\]
Y la desigualdad triangular de la métrica es la de la norma aplicada a
$x - z = (x-y) + (y-z)$.
\end{proof}
}

\begin{remark}
El recíproco es falso: no toda métrica sobre un espacio vectorial proviene de
una norma.
\onlynotes{%
  Basta observar que una métrica inducida cumple $d(\lambda x, \lambda y) =
  |\lambda|\, d(x,y)$, y que la métrica discreta no.}
\end{remark}
\end{frame}
